\section{Profiling Interface}
\mpitermtitleindex{profiling interface}

\label{sec:prof}
\label{chap:prof}

\subsection{Requirements}
\label{sec:prof:require}

To meet the requirements for the \MPI/ profiling interface, an
implementation of the \MPI/ functions \emph{must}
\begin{enumerate}
\item provide a mechanism through which all of the \MPI/ defined functions,
  except those allowed as macros (See Section~\ref{sec:macros}), may be
  accessed with a name shift.  This requires, in C and Fortran, an
  alternate entry point name, with the prefix \mpifuncskip{PMPI\_} for each \mpi/
  \mpitermdefindex{PMPI\_} function in each provided language binding and
  language support method.  For routines implemented as macros, it is still
  required that the \mpifuncskip{PMPI\_} version be supplied and work as
  expected, but it is not possible to replace the
  \mpifuncskip{MPI\_} version with a user-defined version at link time.
 
  For Fortran, the different support methods cause several specific
  procedure names.  Therefore, several profiling routines (with these
  specific procedure names) are needed for each Fortran \MPI/ routine, as
  described in \sectionref{sec:f90:linker-names}.

\item ensure that those \MPI/ functions that are not replaced may still be
  linked into an executable image without causing name clashes.

\item document the implementation of different language bindings of the
  \MPI/ interface if they are layered on top of each other, so that the
  profiler developer knows whether to implement the profile interface
  for each binding, or to economize by implementing it only for the lowest
  level routines.

\item where the implementation of different language bindings is done
  through a layered approach (e.g., the Fortran binding is a set of
  ``wrapper'' functions that  call the C implementation), ensure that these
  wrapper functions are separable from the rest of the library.

  This separability is necessary to allow a separate profiling library to
  be correctly implemented, since (at least with Unix linker semantics) the
  profiling library must contain these wrapper functions if it is to
  perform as expected.  This requirement allows the person who builds the
  profiling library to extract these functions from the original \MPI/
  library and add them into the profiling library without bringing along
  any other unnecessary code.

\item provide a no-op routine \mpifunc{MPI\_PCONTROL} in the \MPI/ library.
\end{enumerate}

\subsection{Discussion}

The objective of the \MPI/ profiling interface is to ensure that it is
relatively easy for authors of profiling (and other similar) tools to
interface their codes to \MPI/ implementations on different machines.

Since \MPI/ is a machine independent standard with many different
implementations, it is unreasonable to expect that the authors of profiling
tools for \MPI/ will have access to the source code that  implements \MPI/
on any particular machine.  It is therefore necessary to provide a
mechanism by which the implementors of such tools can collect whatever
performance information they wish \emph{without} access to the underlying
implementation.

We believe that having such an interface is important if \MPI/ is to be
attractive to end users, since the availability of many different tools
will be a significant factor in attracting users to the \MPI/ standard.

The profiling interface is just that, an interface.  It says \emph{nothing}
about the way in which it is used.  There is therefore no attempt to lay
down what information is collected through the interface, or how the
collected information is saved, filtered, or displayed.

While the initial impetus for the development of this interface arose from
the desire to permit the implementation of profiling tools, it is clear
that an interface like that specified may also prove useful for other
purposes, such as ``internetworking'' multiple \MPI/ implementations.
Since all that is defined is an interface, there is no objection to it
being used wherever it is useful.

As the issues being addressed here are intimately tied up with the way in
which executable images are built, which may differ greatly on different
machines, the examples given below should be treated solely as one way of
implementing the objective of the \MPI/ profiling interface.  The actual
requirements made of an implementation are those detailed in the
Requirements section above, the whole of the rest of this section is only
present as justification and discussion of the logic for those
requirements.

Examples~\ref{exa:prof:measurement-wrapper}, \ref{exa:prof:weak-symbols}, \ref{exa:prof:macros}, and~\ref{exa:prof:linking} show ways in which an implementation could be
constructed to meet the requirements on a Unix system (there may be others that would be equally valid).

\subsection{Logic of the Design}

Provided that an \MPI/ implementation meets the requirements above, it is
possible for the implementor of the profiling system to intercept the \MPI/
calls that  are made by the user program.  The profiling system implementor can then collect any required
information before calling the underlying \MPI/ implementation
(through its name shifted entry points) to achieve the desired effects.


\begin{example}
\label{exa:prof:measurement-wrapper}
\exindex{Profiling interface!measurement wrapper}%
\Exindex{C}{Measurement wrapper}{MPI\_Type\_size,MPI\_Send,MPI\_Wtime}%
A wrapper to accumulate the total amount of data sent by the
\mpifunc{MPI\_SEND} function, along with the total elapsed time spent in
the function.

%%HEADER
%%LANG: C
%%ENDHEADER
\begin{lstlisting}[language={[MPI]C}]
static int totalBytes = 0;
static double totalTime = 0.0;

int MPI_Send(const void* buffer, int count, MPI_Datatype datatype,
             int dest, int tag, MPI_Comm comm)
{
    double tstart = MPI_Wtime();       /* Pass on all arguments */
    int size;
    int result    = PMPI_Send(buffer, count, datatype, dest, tag, comm);

    totalTime    += MPI_Wtime() - tstart;  /* Compute time */

    MPI_Type_size(datatype, &size);        /* and size */
    totalBytes += count*size;

    return result;
}
\end{lstlisting}
\end{example}

\subsection{\texorpdfstring{\MPI/}{MPI} Library Implementation}

If the \MPI/ library is implemented in C on a Unix system, then there are
various options, including the two presented here, for supporting the
name-shift requirement.  The choice between these two options depends
partly on whether the linker and compiler support weak symbols.

If the compiler and linker support weak external symbols, then only 
a single library is required as the following example shows: 

\begin{example}
\label{exa:prof:weak-symbols}
\Exindex[Profiling interface, implementation using weak symbols]{C}{Profiling interface!implementation using weak symbols}{}
Library implementation using weak symbols.
%%HEADER
%%LANG: C
%%SUBST:/\* Useful content \*/:return 1;
%%SUBST:/\* appropriate args \*/:void
%%ENDHEADER
\begin{lstlisting}[language={[MPI]C}]
#pragma weak MPI_Example = PMPI_Example

int PMPI_Example(/* appropriate args */)
{
    /* Useful content */
}
\end{lstlisting}
\end{example}

The effect of this \code{\#pragma} is to define the external symbol
\code{MPI\_Example} as a weak definition.  This means that the linker will
not complain if there is another definition of the symbol (for instance in
the profiling library); however if no other definition exists, then the
linker will use the weak definition.

In the absence of weak symbols then one possible solution would be to use
the C macro preprocessor as the following example shows:

\begin{example}
\label{exa:prof:macros}
\Exindex[Profiling interface, implementation using the C macro preprocessor]{C}{Profiling interface!implementation using the C macro preprocessor}{}
Library implementation using C pre-processor macros.
%%HEADER
%%LANG: C
%%ENDHEADER
\begin{lstlisting}[language=C]
#ifdef PROFILELIB
#    ifdef __STDC__
#        define FUNCTION(name) P##name
#    else
#        define FUNCTION(name) P/**/name
#    endif
#else
#    define FUNCTION(name) name
#endif
\end{lstlisting}

Each of the user visible functions in the library would then be declared
thus

%%HEADER
%%LANG: C
%%TOP:#define FUNCTION(name) name
%%SUBST:/\* Useful content \*/:return 1;
%%SUBST:/\* appropriate args \*/:void
%%ENDHEADER
\begin{lstlisting}[language={[MPI]C}]
int FUNCTION(MPI_Example)(/* appropriate args */)
{
    /* Useful content */
}
\end{lstlisting}

The same source file can then be compiled to produce both versions of the
library, depending on the state of the \code{PROFILELIB} macro symbol.
\end{example}

It is required that the standard \MPI/ library be built in such a way that
the inclusion of \MPI/ functions can be achieved one at a time.  This may mean that each external function must reside in its own compilation unit. This is
necessary so that the author of the profiling library need only define
those \MPI/ functions that need to be intercepted, references to any
others being fulfilled by the normal \MPI/ library.

\begin{example}
\label{exa:prof:linking}
  \exindex{Linking libraries when using the profiling interface}

	The following example shows a potential link step when using the profiling
	interface.

%%HEADER
%%SKIP
%%ENDHEADER
\begin{lstlisting}
% cc ... -lmyprof -lpmpi -lmpi
\end{lstlisting}

Here \code{libmyprof.a} contains the profiler functions that  intercept
some of the \MPI/ functions, \code{libpmpi.a} contains the ``name shifted''
\MPI/ functions, and \code{libmpi.a} contains the normal definitions of the
\MPI/ functions.

\end{example}

\subsection{Complications}
\label{sec:prof:complications}

\subsubsection{Multiple Counting}

Since parts of the \MPI/ library may themselves be implemented using more
basic \MPI/ functions (e.g., a portable implementation of the collective
operations implemented using point-to-point communications), there is
potential for profiling functions to be called from within an \MPI/
function that  was called from a profiling function.  This could lead to
``double counting'' of the time spent in the inner routine.  Since this
effect could actually be useful under some circumstances (e.g., it might
allow one to answer the question ``How much time is spent in the
point-to-point routines when they are called from collective functions?''),
we have
decided not to enforce any restrictions on the author of the \MPI/ library
that  would overcome this.  Therefore the author of the profiling library
should be aware of this problem, and guard against it.  In a
single-threaded world this is easily achieved through use of a static
variable in the profiling code that  remembers if you are already inside a
profiling routine.  It becomes more complex in a multithreaded environment
(as does the meaning of the times recorded).

\subsubsection{Linker Oddities}

The Unix linker traditionally operates in one pass: the effect of this is
that functions from libraries are only included in the image if they are
needed at the time the library is scanned.  When combined with weak
symbols, or multiple definitions of the same function, this can cause odd
(and unexpected) effects.

Consider, for instance, an implementation of \MPI/ in which the Fortran
binding is achieved by using wrapper functions on top of the C
implementation.  The author of the profile library then assumes that it is
reasonable only to provide profile functions for the C binding, since
Fortran will eventually call these, and the cost of the wrappers is assumed
to be small.  However, if the wrapper functions are not in the profiling
library, then none of the profiled entry points will be undefined when the
profiling library is called.  Therefore none of the profiling code will be
included in the image.  When the standard \MPI/ library is scanned, the
Fortran wrappers will be resolved, and will also pull in the base versions
of the \MPI/ functions.  The overall effect is that the code will link
successfully, but will not be profiled.

To overcome this we must ensure that the Fortran wrapper functions are
included in the profiling version of the library.  We ensure that this is
possible by requiring that these be separable from the rest of the base
\MPI/ library.  This allows them to be copied out of the base library and
into the profiling one using a tool such as \code{ar}.

\subsubsection{Fortran Support Methods}

The different Fortran support methods and possible options for the support
of subarrays (depending on whether the compiler can support \code{TYPE(*),
DIMENSION(..)} choice buffers) imply different specific procedure names for
the same Fortran \MPI/ routine.  The rules and implications for the
profiling interface are described in \sectionref{sec:f90:linker-names}.
 
\subsection{Multiple Levels of Interception}

The scheme given here does not directly support the nesting of profiling
functions, since it provides only a single alternative name for each \MPI/
function.  Consideration was given to an implementation that  would allow
multiple levels of call interception, however we were unable to construct
an implementation of this that  did not have the following disadvantages
\begin{itemize}
\item
  assuming a particular implementation language, and
\item
  imposing a run time cost even when no profiling was taking place.
\end{itemize}
Since one of the objectives of \MPI/ is to permit efficient, low latency
implementations, and it is not the business of a standard to require a
particular implementation language, we decided to accept the scheme
outlined above.

Note, however, that it is possible to use the scheme above to implement a
multi-level system, since the function called by the user may call many
different profiling functions before calling the underlying \MPI/ function.
This capability has been demonstrated in the P$^N$MPI tool
infrastructure~\cite{schulz+:sc07}.

\subsection{Miscellaneous Control of Profiling}

There is a clear requirement for the user code to be able to control the
profiler dynamically at run time.  This capability is normally used for (at
least) the purposes of
\begin{itemize}
\item Enabling and disabling profiling depending on the state of the
  calculation.
\item Flushing trace buffers at noncritical points in the calculation.
\item Adding user events to a trace file.
\end{itemize}
These requirements are met by use of \mpifunc{MPI\_PCONTROL}.

\begin{mpi-binding}
    function_name("MPI_Pcontrol")
    parameter("level", "PROFILE_LEVEL", desc=r"Profiling level", constant=True)
    parameter("varargs", "VARARGS")
    no_ierror()
\end{mpi-binding}

\MPI/ libraries themselves make no use of this routine, and simply return
immediately to the user code.  However the presence of calls to this
routine allows a profiling package to be explicitly called by the user.

Since \MPI/ has no control of the implementation of the profiling code, we
are unable to specify precisely the semantics that  will be provided by
calls to \mpifunc{MPI\_PCONTROL}.  This vagueness extends to the number of
arguments to the function, and their datatypes.

However to provide some level of portability of user codes to different
profiling libraries, we request the following meanings for certain values
of level.
\begin{itemize}
\item{\code{level=0}}
  Profiling is disabled.
\item{\code{level=1}}
  Profiling is enabled at a normal default level of detail.
\item{\code{level=2}}
  Profile buffers are flushed, which may be a no-op in some profilers.
\item{All other values of \code{level}}
  have profile library defined effects and additional arguments.
\end{itemize}

We also request that the default state after \MPI/ has been
initialized is for profiling to be enabled at the normal default level.  (i.e.,
as if \mpifunc{MPI\_PCONTROL} had just been called with the argument 1).
This allows users to link with a profiling library and to obtain profile
output without having to modify their source code at all.

The provision of \mpifunc{MPI\_PCONTROL} as a no-op in the standard \MPI/
library supports the collection of more detailed profiling information with
source code that can still link against the standard \MPI/ library.
